% Run GCT_KRON_20261120; GPT-6 Astra Pro; effort not exposed.
\section{Three permanent copies at a positive rectangular Kronecker coefficient}

\subsection*{A-033. A uniform analytic multiplicity obstruction after doubling}
For \(n\geq3\), use the A-030 notation
\[
 R=n^2,\quad D=R+1,\quad q=R(n-2),\quad E=n+q=n(n-1)^2,
 \qquad\lambda_n=(E+1,3,2^{R-2}).
\]
Then
\[
 \boxed{\quad
 q^n_{2\lambda_n}(2D[n])\geq3
   >g(2\lambda_n,((2D)^n),((2D)^n))=1
   \geq K_n(2\lambda_n).
 \quad}                                                       \tag{A-033.1}
\]
Thus the obstruction persists at a weight with a positive
ordinary rectangular Kronecker coefficient. The argument works
for every \(n\geq3\) and uses explicit polynomial identities,
without a finite-rank computation.
It concerns the unpadded case \(n=m\).
At \(n=4\) the determinant multiplicity is known to be positive:
\[
 q^4_{(74,6,4^{14})}(34[4])\geq3
       >K_4(74,6,4^{14})
       =\operatorname{sk}((74,6,4^{14}),(34^4))
       =g((74,6,4^{14}),(34^4),(34^4))=1.                       \tag{A-033.2}
\]
The determinant-side equality in (A-033.2) is the known-invariant
specialization of O-028, recalled below. No exact value of the
permanent multiplicity is asserted. A-034 below subsequently
determines \(K_n(s\lambda_n)\) for every \(n\geq3\) and every
stretch \(s\), and strengthens the actual obstruction to that
whole family.

\paragraph{Three universal highest-weight polynomials.}
Fix an ordered basis of \(W=\operatorname{Mat}_n(\mathbb C)\).
For \(F\in\operatorname{Sym}^nW^*\), let
\[
                 H(F)=\det\operatorname{Hess}_{R}(F)
\]
be the Hessian determinant in all \(R\) variables.
It has degree \(R\) in the coefficients of \(F\) and order
\(q\) in the variables. In particular it is not the Hessian
of the following binary restriction.
Define normalized coefficients by
\[
 \begin{split}
 F(xe_1+ye_2)&=\sum_{i=0}^n\binom ni F_i\,x^{n-i}y^i,\\
 H(F)(xe_1+ye_2)&=\sum_{i=0}^{q}\binom qi H_i\,x^{q-i}y^i.
 \end{split}                                                   \tag{A-033.3}
\]
The \(F_i\) are linear coefficient functions and the \(H_i\)
are polynomial coefficient functions of degree \(R\).
Consider
\[
 \begin{split}
 W_1&=(F_0F_2-F_1^2)H_0^2,\\
 W_2&=F_0^2(H_0H_2-H_1^2),\\
 W_3&=(F_0H_1-F_1H_0)^2.
 \end{split}                                                   \tag{A-033.4}
\]
Each has coefficient degree \(2+2R=2D\).
We verify their highest-weight character on the full group
\(\operatorname{GL}(W)\), rather than just on the binary subgroup.

Let \(g\) be upper triangular, with
\(ge_1=a e_1\) and \(ge_2=b e_1+c e_2\).
For the coordinate-ring action evaluate \(W_i\) at \(F\circ g\).
The first three normalized binary coefficients of \(F\circ g\) are
\[
 a^nF_0,\qquad
 a^{n-1}(bF_0+cF_1),\qquad
 a^{n-2}(b^2F_0+2bcF_1+c^2F_2).
                                                               \tag{A-033.5}
\]
The same formulas hold for \(H_0,H_1,H_2\), replacing
\(n\) by \(q\) and multiplying every coefficient by
\(\det(g)^2\), because
\[
                         H(F\circ g)=\det(g)^2 H(F)\circ g.
\]
Substitution in (A-033.4) gives, for all three polynomials,
\[
                 W_i(F\circ g)=
           \det(g)^4 a^{2E-2}c^2\,W_i(F).                      \tag{A-033.6}
\]
Thus each nonzero \(W_i\) is a highest-weight vector of weight
\[
                 (2E+2,6,4^{R-2})=2\lambda_n.                  \tag{A-033.7}
\]
The determinant character here is essential: omitting it
would identify only the two-row core, not the full
\(\operatorname{GL}_{n^2}\) highest weight.

\paragraph{A one-parameter family inside a fixed special-linear orbit.}
Let \(P=\operatorname{per}_n\), let \(J\) be the all-ones matrix,
and write \(Z=E_{11}\). For \(s\in\mathbb C\), set
\[
 A_s=J+(s-n+2)Z,\qquad
                       M_s(x,y)=xA_s+yZ.                      \tag{A-033.8}
\]
Every entry except \(11\) equals \(x\), while entry \(11\)
is \((s-n+3)x+y\).
The matrices \(A_0,Z\) are linearly independent.
Extend them to the first two columns of some
\(g_0\in\operatorname{SL}(W)\), scaling one remaining column
of an arbitrary completion to make its determinant one.
Then \(A_s=A_0+sZ\) shows that these two columns extend to
\(g_s=g_0u_s\in\operatorname{SL}(W)\), where
\(u_se_1=e_1+s e_2\), \(u_se_2=e_2\), and \(u_s\) fixes
the other basis vectors.
Consequently the binary data below are evaluations of
(A-033.3) at \(P\circ g_s\), all on one fixed
special-linear orbit. The full Hessian has no extra scalar
factor since \(\det(g_s)=1\).

Put
\[
 h=(n-2)^2,\qquad
 \alpha=(n-1)!,\qquad
 \beta=((n-3)!)^{n^2}(n-2)^{4(n-1)}(n-1)^{2n}.
\]
The defining permanent sum and the exact Hessian formula
(A-030.13), with homogeneity, give
\[
 \begin{split}
 P(M_s(x,y))&=\alpha x^{n-1}((s+2)x+y),\\
 H(P)(M_s(x,y))&=\beta x^{q-h}(sx+y)^h.
 \end{split}                                                   \tag{A-033.9}
\]
Here \(\alpha,\beta\ne0\), \(h\geq1\), and \(q>h\).
The parameter in (A-033.8) specifies the \(11\) entry;
in a pencil written instead as \(xJ+(t x+y)E_{11}\), the
corresponding parameter is \(s=t+n-2\).
This fixes the possible one-unit shift in using the
all-ones normalization.

From (A-033.3) and (A-033.9),
\[
 F_0=\alpha(s+2),\quad F_1=\alpha/n,\quad F_2=0,\qquad
 H_0=\beta s^h,\quad H_1=\beta(h/q)s^{h-1}.
\]
For \(h\geq2\),
\(H_2=\beta h(h-1)s^{h-2}/(q(q-1))\);
for \(h=1\), \(H_2=0\).
Thus the three evaluations in (A-033.4) have the common
nonzero polynomial factor \(\alpha^2\beta^2s^{2h-2}\),
and their remaining factors are
\[
 -\frac{s^2}{n^2},\qquad
 \frac{h(h-q)}{q^2(q-1)}(s+2)^2,\qquad
 \left(\left(\frac hq-\frac1n\right)s+\frac{2h}q\right)^2
                 =\frac4{n^4}(s-(n-2))^2.                    \tag{A-033.10}
\]
These formulas are polynomial identities also at \(s=0\).
Every scalar multiplier of the three displayed squares is
nonzero for \(n\geq3\).

The coefficient matrix of
\(s^2,(s+2)^2,(s-(n-2))^2\), in the order \(s^2,s,1\), is
\[
 \begin{pmatrix}
 1&0&0\\
 1&4&4\\
 1&-2(n-2)&(n-2)^2
 \end{pmatrix},
 \qquad\det=4n(n-2)\ne0.                                      \tag{A-033.11}
\]
Hence the three evaluations are linearly independent.
A relation among the restrictions of \(W_1,W_2,W_3\)
to the permanent orbit closure would restrict to a relation
among these three polynomials in \(s\).
Cancellation of the common nonzero polynomial factor in
\(\mathbb C[s]\) and (A-033.11) force all its coefficients
to vanish. This argument does not require the common factor
to be nonzero at every individual parameter value.

We have therefore produced three independent highest-weight
vectors of type \(2\lambda_n\) in the actual permanent
coordinate ring. Its multiplicity is at least three.
A-030.15 gives
\[
 g(2\lambda_n,((2D)^n),((2D)^n))
             =[z^2]\prod_{j=2}^n(1-z^j)^{-1}=1.
 \]
Together with \(K_n\leq\operatorname{sk}\leq g\), this proves
(A-033.1).

\paragraph{Products give a lower bound at every dilation.}
For \(u\geq0\), the same three evaluations prove
\[
 q^n_{2u\lambda_n}(2uD[n])\geq2u+1,\qquad
 q^n_{(2u+1)\lambda_n}((2u+1)D[n])\geq2u+1.                    \tag{A-033.12}
\]
The first inequality includes the constant component at \(u=0\).
Indeed the three factors in (A-033.10) span all polynomials
of degree at most two in \(s\). Products of \(u\) such factors
span all polynomials of degree at most \(2u\), a space of
dimension \(2u+1\). Multiplication by their common nonzero
factor to the \(u\)-th power preserves that dimension.
Every such product of the \(W_i\) is a highest-weight vector
of weight \(2u\lambda_n\), so restriction already has rank at
least \(2u+1\). This proves the even case.

For the odd case, multiply that whole highest-weight space by
the nonzero weight-\(\lambda_n\), degree-\(D\) section from
A-030. The permanent orbit closure is irreducible, and its
coordinate ring is a domain. This multiplication is injective
and proves the stated odd lower bound.
No algebraic independence of the three \(W_i\) is assumed:
the argument uses the dimension of their product span on the
specific one-parameter family.

In particular, when \(n=3\), these bounds give an analytic
proof of
\[
        q^3_{s\lambda_3}(10s[3])
              >g(s\lambda_3,((10s)^3),((10s)^3))
              \geq K_3(s\lambda_3)
              \qquad(s\geq1).                                \tag{A-033.13}
\]
For \(s\geq2\), the ordinary coefficient counts solutions
\(2a+3b=s\), so is at most \(\lfloor s/3\rfloor+1<s\).
The lower bound in (A-033.12) is \(s+1\) for even \(s\),
and \(s\) for odd \(s\). At \(s=1\), A-030 gives \(q=1>g=0\).
This establishes persistence at every positive dilation.
F-028 supplies a stronger numerical lower bound in this
cubic case, by a separate finite dominance certificate.
O-028 shows that the determinant multiplicity in (A-033.13)
is positive and equals the ordinary coefficient for every
sufficiently large \(s\), with no numerical threshold claimed.

\paragraph{The explicit positive determinant side at size four.}
O-028 uses the known invariant computation
\[
                            P_{4,16}(\det_4)\ne0
\]
reported by B\"urgisser--Ikenmeyer,
\emph{Fundamental invariants of orbit closures}.
In arXiv:1511.02927v2, see the paragraph following Proposition 3.28,
printed p.~18, together with Theorem 3.15, p.~13.
This gives an actual determinant-coordinate invariant of degree
\(16\), with highest weight \(4\mathbf1_{16}\).
That published nonvanishing is not recomputed here.
Multiplication by this invariant injects the nonzero two-row
space of weight \((70,2)\), degree \(18\), into the weight
\[
                    (70,2)+4\mathbf1_{16}
                             =(74,6,4^{14}),
\]
in degree \(34\). Its two-row multiplicity is one by O-006.
The ordinary upper bound for the shifted weight is also one,
by A-030.15 or O-028's determinant twist.
Thus \(K_4=\operatorname{sk}=g=1\), proving (A-033.2).
The primary source and its source-check record are those of
O-028; the present new input is the independent permanent
lower bound in (A-033.10)--(A-033.11).

\paragraph{Geometric dictionary and scope.}
The functions in (A-033.4) are actual homogeneous polynomial
functions on the form space. The Hessian is a covariant with
a determinant-line factor, and the normalized binary
coefficient operations are polynomial over \(\mathbb Q\).
Their upper-triangular characters identify complete
highest-weight multiplicity vectors. The exact dictionary is
SGA 3, I, Proposition 4.7.2, pp.~34--35, for polynomial
comodule maps, and Proposition 4.7.3, pp.~35--36, for their
torus gradings; SGA 6, VI, Theorem 3.3, pp.~392--393,
for the Schur and exterior operations; and Macdonald,
I, Appendix A(5.3)--(5.6), pp.~155--156, for the
characteristic-zero Schur decomposition, as used in F-016.
Restriction to the explicit special-linear orbit family proves
linear independence in the quotient coordinate ring.
Products then use its domain property.
Every conclusion concerns actual multiplicity spaces, rather
than positivity of virtual classes.
All weights \(s\lambda_n\) have length \(n^2\) for \(s>0\);
their internally padded permanent multiplicity is therefore
zero when \(m<n\). No transfer to a large-padding range is made.
